En algunes missions espacials s'han utilitzat motors iònics. En aquests motors es produeixen ions positius que s'envien a una cambra on un camp elèctric constant els impulsa. El motor expulsa ions positius a gran velocitat i la nau adquireix impuls en sentit contrari. Considereu un motor iònic en què ions Xe$^{+}$, inicialment en un estat de repòs, s'acceleren entre dues plaques separades 10~cm fins a adquirir una velocitat de $3,0\times 10^{5}\ \mathrm{m/s}$.

\figura[0.4]{fig1.png}

\begin{apartat}{1}
Calculeu l'acceleració dels ions i el camp elèctric (que podeu considerar constant) a la cambra d'acceleració.
\end{apartat}

\begin{apartat}{1}
Calculeu la diferència de potencial entre les dues plaques amb les dades de la figura. Indiqueu també el valor que hauria de tenir aquesta diferència de potencial si les dues plaques estiguessin separades només 6~cm per a aconseguir la mateixa velocitat de sortida dels ions.
\end{apartat}

\dades{%
  $Q(\text{ions Xe}^{+}) = +1,60\times 10^{-19}\ \mathrm{C}$\\
  $m(\text{ions Xe}^{+}) = 132\ \mathrm{u}$\\
  $1\ \mathrm{u} = 1,66\times 10^{-27}\ \mathrm{kg}$}
